\section{Generador de Confirmación de Enfermero}

Este componente modela la intervención reactiva y estocástica del personal de enfermería ante las alertas del sistema. Se adopta un diseño estrictamente \textbf{reactivo-estocástico}. El componente permanece ocioso en un estado pasivo hasta la recepción de un estímulo del sistema clínico. Ante la ocurrencia de una alarma, el modelo transiciona a una fase activa donde programa internamente un evento de confirmación médica retrasado por una latencia aleatoria uniforme.

\subsection{Definición formal}

\paragraph{Puertos y estructura del estado.}
\begin{align*}
X &= \{\text{alarma}\} \times \{0\} \\[4pt]
Y &= \{\mathit{confirmacionEnfermero}\} \times \{0\} \\[4pt]
S &= \{\text{idle}, \text{esperandoConfirmacion}\} \times (\mathbb{R}_0^+ \cup \{\infty\})
\end{align*}

La tupla de estado se formaliza como $s = (\mathit{fase}, \sigma)$, donde:
\begin{itemize}
    \item $\mathit{fase} \in \{\text{idle}, \text{esperandoConfirmacion}\}$: Indica si el componente se encuentra en reposo pasivo o si posee una acción de confirmación.
    \item $\sigma \in \mathbb{R}_0^+ \cup \{\infty\}$: Tiempo remanente hasta la emisión física del evento de confirmación.
\end{itemize}

\paragraph{Estado inicial.} 
\[
s_0 = (\text{idle}, \infty)
\]

\paragraph{Transición externa ($\dext$).}
\[
\dext \bigl( (\mathit{fase}, \sigma), e, s_x \bigr) =
\begin{cases}
(\text{esperandoConfirmacion}, T_{\mathrm{resp}}) & \text{si } s_x = (\text{alarma}, 0) \text{ y } \mathit{fase} = \text{idle} \\[6pt]
(\mathit{fase}, \sigma - e) & \text{si } s_x = (\text{alarma}, 0) \text{ y } \mathit{fase} = \text{esperandoConfirmacion}
\end{cases}
\]

La distribución temporal utilizada es estocástica dentro de un intervalo uniforme continuo:
\[
T_{\mathrm{resp}} \sim U(t_{\mathrm{min}}, t_{\mathrm{max}}) = \text{rand}(5, 75)
\]

\paragraph{Origen de la estimación.}
La selección de la ventana operativa $[5, 75]$ segundos segmenta la simulación en tres zonas de comportamiento dinámico:
\begin{itemize}
    \item \textbf{Zona de Respuesta Oportuna ($5 \le T_{\mathrm{resp}} \le 30\ \mathrm{s}$):} Evalúa el régimen operativo normal donde el factor humano reacciona con eficiencia.
    \item \textbf{Zona de Reincidencia de Alarmas ($30 < T_{\mathrm{resp}} \le 60\ \mathrm{s}$):} Modela retrasos moderados por parte del personal médico. 
    \item \textbf{Zona de Mitigación Autónoma ($60 < T_{\mathrm{resp}} \le 75\ \mathrm{s}$):} Supera intencionalmente el umbral de seguridad de 60 segundos. 
\end{itemize}

\paragraph{Transición interna ($\dint$).}
\[
\dint(\mathit{fase}, \sigma) = (\text{idle}, \infty)
\]

\paragraph{Función de salida ($\lambda$).}
\[
\lambda(\mathit{fase}, \sigma) =
\begin{cases}
(\mathit{confirmacionEnfermero}, 0) & \text{si } \mathit{fase} = \text{esperandoConfirmacion} \\[4pt]
\varnothing & \text{en otro caso}
\end{cases}
\]

\paragraph{Avance temporal ($\ta$).}
\[
\ta(\mathit{fase}, \sigma) = \sigma
\]